% -*- coding: utf-8 -*-
% XeLaTeX 编译
\documentclass[a3paper, landscape, 11pt, twocolumn]{article}

% 引入宏包
\usepackage{fontspec}
\usepackage[UTF8]{ctex}
\usepackage{geometry}
\usepackage{listings}
\usepackage{xcolor}
\usepackage{enumitem}
\usepackage{tabularx}
\usepackage{makecell}
\usepackage{fancyhdr}
\usepackage{amssymb}
\usepackage{lastpage}
\usepackage{titlesec}

% 设置页面边距 (A3横向双栏)
\geometry{left=2cm, right=2cm, top=2.5cm, bottom=2.5cm, columnsep=1.8cm}

% 颜色定义
\definecolor{codebg}{RGB}{245,245,245}
\definecolor{keywordcolor}{RGB}{0,0,150}
\definecolor{commentcolor}{RGB}{0,128,0}
\definecolor{stringcolor}{RGB}{163,21,21}

% 代码样式（备用）
\lstdefinestyle{JavaExam}{
    language=Java,
    basicstyle=\ttfamily\small,
    backgroundcolor=\color{codebg},
    keywordstyle=\color{keywordcolor}\bfseries,
    commentstyle=\color{commentcolor},
    stringstyle=\color{stringcolor},
    showstringspaces=false,
    breaklines=true,
    frame=single,
    rulecolor=\color{black!50},
    tabsize=4,
    xleftmargin=1em,
    xrightmargin=1em,
    aboveskip=1em,
    belowskip=1em
}

% 页眉页脚设置
\pagestyle{fancy}
\fancyhf{}
\fancyhead[C]{\Large\bfseries 《面向对象程序设计》简答题专项练习}
\fancyfoot[C]{第 \thepage\ 页 共 \pageref{LastPage}\ 页}
\fancyfoot[R]{得分：\underline{\hspace{2cm}}}
\fancyfoot[L]{题号：\underline{\hspace{1cm}}}
\renewcommand{\headrulewidth}{0.4pt}

% 简答题环境
\newcounter{shortq}
\setcounter{shortq}{0}
\newcommand{\sq}{%
    \refstepcounter{shortq}%
    \noindent\textbf{\theshortq、}\hspace{0.5em}%
}

% 答题空白行
\newcommand{\answerblank}{\vspace{3cm}}

\begin{document}

% 试卷头部
\twocolumn[{
    \begin{center}
        \vspace*{-1cm}
        {\huge\bfseries 《面向对象程序设计》简答题专项练习}\\[0.8em]
        {\large 考试时间：120 分钟 \hspace{2em} 总分：150 分}\\[1em]
        {\large 姓名：\underline{\hspace{4cm}} \hspace{2em} 学号：\underline{\hspace{5cm}} \hspace{2em} 班级：\underline{\hspace{4cm}}}
        \vspace{0.8em}
        \hrule
        \vspace{1em}
    \end{center}
}]

\section*{【简答题】共 30 小题，每小题 5 分，共 150 分。}
\textit{（请将答案写在题目下方的空白处，要求要点明确、条理清晰）}

% 题目开始
\sq 简述 Java 语言的主要特点。 \answerblank

\sq 解释 JDK、JRE 和 JVM 的区别与联系。 \answerblank

\sq 什么是标识符？Java 对标识符的命名规则有哪些？ \answerblank

\sq Java 的基本数据类型有哪些？并说明每种类型的占用字节数。 \answerblank

\sq 简述自动类型转换和强制类型转换的区别，并举例说明。 \answerblank

\sq 解释 `==` 和 `equals()` 方法的区别。 \answerblank

\sq 数组的 `length` 属性和字符串的 `length()` 方法有何不同？ \answerblank

\sq 什么是类？什么是对象？二者之间的关系是什么？ \answerblank

\sq 简述构造方法的特点和作用。构造方法可以重载吗？ \answerblank

\sq `this` 关键字有哪些用途？请举例说明。 \answerblank

\sq 什么是方法重载？方法重载的规则是什么？ \answerblank

\sq 什么是方法重写？方法重写需要遵循哪些规则？ \answerblank

\sq 简述类变量（静态变量）和实例变量的区别。 \answerblank

\sq 静态方法和实例方法有何不同？静态方法能否访问实例变量？为什么？ \answerblank

\sq `final` 关键字可以修饰哪些地方？分别起什么作用？ \answerblank

\sq 什么是访问控制修饰符？Java 中的访问控制修饰符有哪些？各自的作用范围是什么？ \answerblank

\sq 简述封装的概念及其在 Java 中的实现方式。 \answerblank

\sq 什么是继承？Java 中如何实现继承？子类能从父类继承哪些成员？ \answerblank

\sq `super` 关键字的作用是什么？请举例说明。 \answerblank

\sq 什么是多态？Java 中多态的实现方式有哪些？ \answerblank

\sq 什么是抽象类？抽象类有哪些特点？ \answerblank

\sq 什么是接口？接口和抽象类的区别是什么？ \answerblank

\sq Java 中如何实现多重继承？为什么？ \answerblank

\sq 什么是异常？Java 的异常处理机制中有哪些关键字？分别说明其作用。 \answerblank

\sq 简述受检查异常（checked exception）和运行时异常（unchecked exception）的区别。 \answerblank

\sq 自定义异常类应如何定义？请写出基本步骤。 \answerblank

\sq `String` 和 `StringBuilder`、`StringBuffer` 的区别是什么？ \answerblank

\sq 列举 Java 集合框架中常用的 List、Set、Map 接口的实现类，并简述它们的特点。 \answerblank

\sq 什么是泛型？使用泛型有什么好处？ \answerblank

\sq 简述线程的生命周期，并说明线程创建的方式。 \answerblank

\sq 什么是线程同步？为什么要进行线程同步？synchronized 关键字可以怎么使用？ \answerblank

\sq 简述 Java 中 I/O 流的分类，并举例说明字节流和字符流的区别。 \answerblank

\sq 什么是序列化？如何实现对象的序列化和反序列化？ \answerblank

\sq 简述 JDBC 操作数据库的基本步骤。 \answerblank

% 说明：共30题，但序号会自动递增。为凑齐30题，以上已列出30题（从1到30）。下面留出空白页。
\newpage
% 确保有30题，序号从1到30已覆盖。

\begin{center}
    \textbf{—— 试卷结束，请仔细检查 ——}
\end{center}

\end{document}